\documentclass[12pt]{article}
\usepackage[T1]{fontenc}
\usepackage[utf8]{inputenc}
\usepackage{lmodern}
\usepackage{textcomp}
\usepackage{enumitem}
\usepackage{geometry}
\geometry{margin=1in}
\begin{document}
\begin{center}
Answer Key - Mathematics - Undergraduate Level Assessment 15 \
\small (Answers are ordered exactly as in the question paper.)
\end{center}

\section*{Multiple Choice}
\begin{enumerate}[label=\arabic*.]
\item \textbf{Answer:} A
\\ \textit{Options:} A) (x - 2)(x + 2)(x - 1)(x + 1); B) (x+$1)^4$; C) (x-1)(x+$1)^3$; D) (x-$1)^3(x$+1)
\\ \textit{Note:} The polynomial $x^4$ + 4 can be factored into the linear factors (x - 2)(x + 2)(x - 1)(x + 1) in Z\_5[x] because, when evaluated in the finite field Z\_5, the roots are found to be 2, -2 (which is 3 in Z\_5), 1, and -1 (which is 4 in Z\_5), allowing the polynomial to be expressed as a product of its linear factors corresponding to these roots.
\item \textbf{Answer:} B
\\ \textit{Options:} A) 4; B) 6; C) 8; D) 9
\\ \textit{Note:} The correct answer is 6 because the number of structurally distinct Abelian groups of a given order can be determined by the partitions of the exponents in the prime factorization of that order, and for the order 72, which factors as \(2^3 \times 3^2\), there are 6 distinct partitions corresponding to the number of different Abelian group structures.
\item \textbf{Answer:} C
\\ \textit{Options:} A) 1; B) 2; C) 3; D) 4
\\ \textit{Note:} There are exactly three nonisomorphic trees with 5 vertices, which can be classified based on their structure: one linear tree (a path), one tree with one vertex connected to three others (a star), and one tree with two vertices connected to two others in a more branched configuration.
\item \textbf{Answer:} C
\\ \textit{Options:} A) 1/2; B) 7/12; C) 5/8; D) 2/3
\\ \textit{Note:} The probability that x \textless{} y can be calculated by determining the area of the region where x is less than y within the rectangle formed by the intervals [0, 3] for x and [0, 4] for y, resulting in an area of 5/8 of the total area of the rectangle, which is 12.
\item \textbf{Answer:} B
\\ \textit{Options:} A) (0, 0, 1); B) (3/7, 3/14, 9/14); C) (7/15, 8/15, 1/15); D) (5/6, 1/3, 1/3)
\\ \textit{Note:} The coordinates (3/7, 3/14, 9/14) represent the point on the plane 2 x + y + 3 z = 3 that minimizes the distance to the origin, which can be determined using the method of Lagrange multipliers or by substituting the plane equation into the distance formula and optimizing it.
\end{enumerate}

\section*{True / False}
\begin{enumerate}[label=\arabic*.]
\setcounter{enumi}{5}
\item \textbf{Answer:} True
\\ \textit{Options:} A) True; B) False
\\ \textit{Note:} The statement is true because, by definition of a normal subgroup, for any element \( a \) in \( G \), the left coset \( aH \) is equal to the right coset \( Ha \), as normal subgroups satisfy the condition \( aH = Ha \) for all \( a \in G \).
\item \textbf{Answer:} True
\\ \textit{Options:} A) True; B) False
\\ \textit{Note:} In an abelian group, the operation is commutative, allowing us to rearrange the terms in the expression (g o $h)^2$ = g o h o g o h to equal $g^2$ o $h^2$.
\item \textbf{Answer:} True
\\ \textit{Options:} A) True; B) False
\\ \textit{Note:} The statement is true because the subgroup 4 Z consists of all multiples of 4, and in the group 2 Z, the cosets of 4 Z are precisely 4 Z and 2 + 4 Z, representing the equivalence classes of integers modulo 4.
\item \textbf{Answer:} False
\\ \textit{Options:} A) True; B) False
\\ \textit{Note:} The statement is false because, when calculating f(8) using the recursion f(2 n) = $n^2$ + f[2(n - 1)], the derived value does not equal 36, but instead yields a different result based on the function's defined relationships.
\item \textbf{Answer:} True
\\ \textit{Options:} A) True; B) False
\\ \textit{Note:} The statement is true because the two abelian groups of order 4, which are the cyclic group of order 4 (Z\_4) and the direct product of two cyclic groups of order 2 (Z\_2 × Z\_2), are the only distinct groups that can be formed under the classification of finite abelian groups, as determined by the Fundamental Theorem of Finitely Generated Abelian Groups.
\end{enumerate}

\section*{Long Form Response}
\begin{enumerate}[label=\arabic*.]
\setcounter{enumi}{10}
\item \textbf{Answer:} The triangle is shown below: [asy] pair A,B,C; A = (0,0); B = (5,0); C = (5,10); draw(A--B --C--A); draw(rightanglemark(C,B,A,16)); label("$A$",A,SW); label("$B$",B,SE); label("$C$",C,N); [/asy] We have $\sin A = \frac{BC}{AC}$ and $\cos A = \frac{AB}{AC}$, so $\sin A = 2\cos A$ gives us $\frac{BC}{AC} = 2\cdot\frac{AB}{AC}$. Multiplying both sides by $AC$ gives $BC = 2 AB$, so $\frac{BC}{AB} = 2$. Finally, we have $\tan A = \frac{BC}{AB} = \boxed{2}$. We also could have noted that $\tan A = \frac{\sin A}{\cos A} = \frac{2\cos A}{\cos A } =\boxed{2}$.
\\ \textit{Note:} correct because it correctly uses the definitions of sine and cosine in a right triangle to establish a relationship between the sides, leading to the conclusion that the tangent of angle A, defined as the ratio of the opposite side to the adjacent side, equals 2.
\item \textbf{Answer:} An eigenvector \( v \) of an invertible matrix \( A \) satisfies the equation \( A v = \lambda v \), where \( \lambda \) is the corresponding eigenvalue. This property extends to the matrices \( 2 A \), \( A^2 \), and \( A^{-1} \). For \( 2 A \), we have \( 2 Av = 2\lambda v \), so \( v \) is an eigenvector with eigenvalue \( 2\lambda \). For \( A^2 \), applying \( A \) again gives \( A^2 v = A(\lambda v) = \lambda(Av) = \lambda^2 v \), indicating \( v \) is an eigenvector with eigenvalue \( \lambda^2 \). For the inverse \( A^{-1} \), since \( A \) is invertible, we can rearrange the eigenvector equation to \( A^{-1}(\lambda v) = v \), leading to \( A^{-1}v = \frac{1}{\lambda}v \), confirming that \( v \) is an eigenvector of \( A^{-1} \) with eigenvalue \( \frac{1}{\lambda} \). Thus, \( v \) remains an eigenvector across these transformations.
\\ \textit{Note:} correct because it demonstrates that an eigenvector \( v \) of an invertible matrix \( A \) maintains its eigenvector status under scalar multiplication and composition with itself, as well as with its inverse, using the defining eigenvalue equation.
\end{enumerate}

\vfill
\noindent\textit{End of Answer Key}
\end{document}